\documentclass[11pt,a4paper]{ctexart}

\usepackage[margin=2.5cm]{geometry}
\usepackage{listings}
\usepackage{xcolor}
\usepackage{booktabs}
\usepackage{hyperref}
\usepackage{fancyhdr}
\usepackage{enumitem}
\usepackage{amsmath}

% 页面设置
\pagestyle{fancy}
\fancyhf{}
\lhead{Missing Semester 2026}
\rhead{Shell 入门练习题解答}
\cfoot{\thepage}
\renewcommand{\headrulewidth}{0.4pt}

% 代码样式
\lstset{
    basicstyle=\ttfamily\small,
    backgroundcolor=\color{gray!5},
    frame=single,
    framesep=4pt,
    rulecolor=\color{gray!40},
    breaklines=true,
    breakatwhitespace=false,
    showstringspaces=false,
    columns=fullflexible,
    keepspaces=true,
    language=bash,
    keywordstyle=\color{blue!70!black},
    commentstyle=\color{green!50!black},
    stringstyle=\color{red!60!black},
    morekeywords={echo, cd, ls, pwd, cat, grep, sed, awk, find, sort, uniq, head, tail, wc, curl, jq, xargs, mkdir, chmod, cp, printf, test, which, man, date, stress, kill, source, export, alias, exit},
    morekeywords=[2]{if, then, else, fi, for, while, do, done, in, set, function},
    keywordstyle=[2]{\color{purple!70!black}},
}

% 内联代码
\newcommand{\code}[1]{\texttt{#1}}

% 题目环境
\newcounter{exercisenum}
\newenvironment{exercise}[1]{
    \refstepcounter{exercisenum}
    \subsection*{题\theexercisenum：#1}
    \addcontentsline{toc}{subsection}{题\theexercisenum：#1}
}{
    \vspace{0.5em}
    \hrule
    \vspace{1em}
}

\title{\textbf{Missing Semester 2026}\\课程概览 + Shell 入门\\练习题解答}
\author{运行环境：macOS + zsh（bash 兼容）}
\date{\today}

\begin{document}

\maketitle

\begin{abstract}
本文档为 \href{https://missing-semester-cn.github.io/2026/course-shell/}{Missing Semester 2026 第一讲「课程概览 + Shell 入门」} 的课后练习题完整解答。所有命令均在终端中实际运行验证，包含命令、输出及原理解释。
\end{abstract}

\tableofcontents
\newpage

\section{基础命令与重定向}

\begin{exercise}{确认 Shell 环境}
\textbf{要求}：本课程要求使用类 Unix Shell（Bash 或 ZSH）。运行 \code{echo \$SHELL} 确认。

\textbf{解答}：
\begin{lstlisting}
echo $SHELL
\end{lstlisting}

\textbf{实际输出}：
\begin{lstlisting}
/bin/zsh
\end{lstlisting}

\textbf{说明}：当前默认 Shell 是 zsh，属于类 Unix Shell，符合课程要求。bash 也可用（版本 3.2.57）。Windows 用户需使用 WSL 或 Linux 虚拟机，不能用 \code{cmd.exe} 或 PowerShell。
\end{exercise}

\begin{exercise}{\code{ls -l} 的输出格式}
\textbf{要求}：\code{ls -l} 的 \code{-l} 选项作用是什么？运行 \code{ls -l /}，每一行最前面的 10 个字符分别代表什么？

\textbf{解答}：
\begin{lstlisting}
ls -l /
\end{lstlisting}

\textbf{实际输出（节选）}：
\begin{lstlisting}
total 10
drwxrwxr-x  42 root  admin  1344 Aug 30 13:55 Applications
drwxr-xr-x@ 39 root  wheel  1248 Apr  6 16:10 bin
lrwxr-xr-x@  1 root  wheel    11 Apr  6 16:10 etc -> private/etc
\end{lstlisting}

\textbf{\code{-l} 选项作用}：以长格式（long format）列出文件，显示文件类型、权限、硬链接数、所有者、所属组、大小、修改时间和文件名。

\textbf{前 10 个字符的含义}：

\begin{table}[h]
\centering
\begin{tabular}{c|l|l}
\toprule
\textbf{位置} & \textbf{含义} & \textbf{示例} \\
\midrule
第1位 & 文件类型 & \code{d}=目录, \code{l}=符号链接, \code{-}=普通文件 \\
第2-4位 & 所有者（user）权限 & \code{rwx} = 读、写、执行 \\
第5-7位 & 所属组（group）权限 & \code{r-x} = 读、执行（无写权限） \\
第8-10位 & 其他用户（others）权限 & \code{r-x} = 读、执行 \\
\bottomrule
\end{tabular}
\end{table}

例如 \code{drwxr-xr-x} 表示：这是一个目录，所有者可读写执行，组用户可读执行，其他用户可读执行。

\noindent\textit{注：第11位可能是 \code{@}（macOS 扩展属性）或 \code{+}（ACL），不属于权限位。}
\end{exercise}

\begin{exercise}{glob 模式匹配}
\textbf{要求}：什么是 glob？测试 \code{ls *.txt}、\code{ls file?.txt}、\code{ls \{a,b,c\}.txt} 等模式。

\textbf{glob 是什么}：glob 是 Shell 提供的\textbf{文件名通配符模式匹配}机制。Shell 在执行命令前，会先将包含通配符的参数展开为匹配的文件名列表，再传给程序。

\begin{table}[h]
\centering
\begin{tabular}{c|l}
\toprule
\textbf{模式} & \textbf{含义} \\
\midrule
\code{*} & 匹配任意长度的任意字符（包括空） \\
\code{?} & 匹配恰好一个任意字符 \\
\code{[abc]} & 匹配 a、b、c 中任意一个字符 \\
\code{[a-z]} & 匹配 a 到 z 范围内任意一个字符 \\
\code{\{a,b,c\}} & 大括号扩展，生成 a、b、c 多个字符串 \\
\bottomrule
\end{tabular}
\end{table}

\textbf{实际测试}：
\begin{lstlisting}
mkdir glob_test && cd glob_test
touch file1.txt file2.txt file3.txt a.txt b.txt c.txt abc.txt ab.txt

ls *.txt        # 匹配所有 .txt 文件
ls file?.txt    # 匹配 file + 任意一个字符 + .txt
ls {a,b,c}.txt  # 大括号扩展为 a.txt b.txt c.txt
ls [ab]*.txt    # 匹配以 a 或 b 开头的 .txt 文件
\end{lstlisting}

\textbf{输出}：
\begin{lstlisting}
=== ls *.txt ===
a.txt  ab.txt  abc.txt  b.txt  c.txt  file1.txt  file2.txt  file3.txt

=== ls file?.txt ===
file1.txt  file2.txt  file3.txt

=== ls {a,b,c}.txt ===
a.txt  b.txt  c.txt

=== ls [ab]*.txt ===
a.txt  ab.txt  abc.txt  b.txt
\end{lstlisting}
\end{exercise}

\begin{exercise}{三种引号的区别}
\textbf{要求}：\code{'单引号'}、\code{"双引号"} 和 \code{\$'ANSI 引号'} 有什么区别？写一条命令，输出同时包含字面量 \code{\$}、\code{!} 和换行符的字符串。

\begin{table}[h]
\centering
\begin{tabular}{c|c|c|c|l}
\toprule
\textbf{引号类型} & \textbf{变量展开} & \textbf{转义序列} & \textbf{历史扩展} & \textbf{适用场景} \\
\midrule
\code{'...'} 单引号 & 否 & 否 & 否 & 纯字面量字符串 \\
\code{"..."} 双引号 & 是 & 否（除少数） & 可能触发 & 需要嵌入变量 \\
\code{\$'...'} ANSI-C & 否 & 是 & 否 & 需要转义序列的字面量 \\
\bottomrule
\end{tabular}
\end{table}

\textbf{实际测试}：
\begin{lstlisting}
# 单引号：完全字面量
echo 'It costs $100! Hello\nWorld'
# 输出: It costs $100! Hello\nWorld

# 双引号：展开变量，\n 不解释
price=100
echo "It costs \$price! Hello\nWorld"
# 输出: It costs $price! Hello\nWorld

# ANSI-C 引号：解释 \n，$ 是字面量
echo $'It costs $100! Hello\nWorld'
# 输出:
# It costs $100! Hello
# World
\end{lstlisting}

\textbf{题目要求的命令}（同时包含字面量 \code{\$}、\code{!} 和换行符）：
\begin{lstlisting}
echo $'价格是 $100!\n这是第二行'
\end{lstlisting}

\textbf{输出}：
\begin{lstlisting}
价格是 $100!
这是第二行
\end{lstlisting}

\noindent\textit{也可以用 \code{printf '价格是 \$100!\textbackslash n这是第二行\textbackslash n'}，printf 默认就解释转义序列。}
\end{exercise}

\begin{exercise}{标准流与重定向}
\textbf{要求}：Shell 有三条标准流：stdin(0)、stdout(1)、stderr(2)。运行 \code{ls /nonexistent /tmp}，把 stdout 和 stderr 分别重定向到两个文件。如何把两者都重定向到同一个文件？

\textbf{三条标准流}：
\begin{itemize}[leftmargin=2em]
    \item \textbf{stdin (0)}：标准输入，默认来自键盘
    \item \textbf{stdout (1)}：标准输出，默认打印到终端
    \item \textbf{stderr (2)}：标准错误，默认也打印到终端（独立于 stdout）
\end{itemize}

\textbf{分别重定向到两个文件}：
\begin{lstlisting}
ls /nonexistent /tmp > stdout.txt 2> stderr.txt
\end{lstlisting}
\begin{itemize}[leftmargin=2em]
    \item \code{> stdout.txt} 等价于 \code{1> stdout.txt}，把标准输出写入 stdout.txt
    \item \code{2> stderr.txt} 把标准错误写入 stderr.txt
\end{itemize}

\textbf{实际输出}：
\begin{lstlisting}
stdout.txt: 包含 /tmp 目录下的文件列表
stderr.txt: 包含 "ls: /nonexistent: No such file or directory"
\end{lstlisting}

\textbf{两者重定向到同一个文件（两种方法）}：

方法1（POSIX 标准，推荐）：
\begin{lstlisting}
ls /nonexistent /tmp > both.txt 2>&1
\end{lstlisting}
\begin{itemize}[leftmargin=2em]
    \item \code{> both.txt} 先把 stdout 指向 both.txt
    \item \code{2>&1} 把 stderr 重定向到当前 stdout 指向的地方（即 both.txt）
    \item \textbf{顺序很重要}：\code{2>&1} 必须在 \code{>} 之后
\end{itemize}

方法2（bash/zsh 简写）：
\begin{lstlisting}
ls /nonexistent /tmp &> both.txt
\end{lstlisting}
\end{exercise}

\begin{exercise}{退出状态与 \code{\&\&} / \code{||}}
\textbf{要求}：\code{\$?} 保存上一条命令的退出状态（0=成功）。\code{\&\&} 仅在前一条成功时执行后一条；\code{||} 仅在前一条失败时执行后一条。写一行命令：仅当 \code{/tmp/mydir} 不存在时才创建它。

\textbf{退出状态}：每个命令执行后返回一个 0-255 的整数，0 表示成功，非 0 表示失败。可用 \code{echo \$?} 查看。

\textbf{逻辑运算符}：
\begin{itemize}[leftmargin=2em]
    \item \code{cmd1 \&\& cmd2}：cmd1 成功（退出码0）才执行 cmd2
    \item \code{cmd1 || cmd2}：cmd1 失败（退出码非0）才执行 cmd2
    \item 两者可组合：\code{cmd1 \&\& cmd2 || cmd3}（类似三元运算符）
\end{itemize}

\textbf{题目要求的命令}：
\begin{lstlisting}
[ -d /tmp/mydir ] || mkdir /tmp/mydir
\end{lstlisting}

\textbf{解释}：
\begin{itemize}[leftmargin=2em]
    \item \code{[ -d /tmp/mydir ]} 测试目录是否存在（\code{test -d} 的简写）
    \item 如果存在，\code{[} 命令返回 0（成功），\code{||} 后面的 \code{mkdir} 不执行
    \item 如果不存在，\code{[} 返回非 0（失败），\code{||} 触发 \code{mkdir /tmp/mydir}
\end{itemize}

\textbf{实际验证}：
\begin{lstlisting}
# 第一次运行：目录不存在，创建成功
[ -d /tmp/mydir ] || mkdir /tmp/mydir
ls -ld /tmp/mydir
# 输出: drwxr-xr-x  2 liuchenxi  wheel  64 Aug 30 16:51 /tmp/mydir

# 第二次运行：目录已存在，不重复创建
[ -d /tmp/mydir ] || mkdir /tmp/mydir
# 无输出，无报错
\end{lstlisting}

\noindent\textit{更简洁的等价写法：\code{mkdir -p /tmp/mydir}（\code{-p} 表示目录已存在时不报错），但题目要求用 \code{\&\&}/\code{||} 逻辑。}
\end{exercise}

\section{脚本编写与权限}

\begin{exercise}{为什么 \code{cd} 必须是 Shell 内建命令？}
\textbf{要求}：为什么 \code{cd} 必须是 Shell 内建命令，而不能是独立程序？（提示：想想子进程能影响和不能影响父进程的哪些状态。）

\textbf{核心原因}：子进程无法修改父进程的工作目录。

\textbf{详细解释}：

\begin{enumerate}[leftmargin=2em]
    \item \textbf{进程独立性}：在 Unix 系统中，每个进程都有自己独立的「当前工作目录」（cwd）。当 Shell 执行一个外部程序时，会通过 \code{fork()} 创建子进程，再用 \code{exec()} 加载程序。子进程继承父进程的 cwd，但这只是一个\textbf{副本}。

    \item \textbf{如果 cd 是独立程序}：
    \begin{itemize}[leftmargin=1.5em]
        \item Shell fork 出子进程
        \item 子进程执行 \code{cd /some/path}，修改的是\textbf{子进程自己}的 cwd
        \item 子进程结束后，Shell（父进程）的 cwd \textbf{完全没有改变}
        \item 结果就是：你执行了 \code{cd /some/path}，但提示符显示的当前目录没变
    \end{itemize}

    \item \textbf{内建命令的优势}：内建命令由 Shell 自身执行，不创建子进程，可以直接修改 Shell 进程自己的 cwd，从而真正改变后续命令的工作目录。
\end{enumerate}

\textbf{验证}：
\begin{lstlisting}
which cd
# 输出: cd: shell built-in command（或 no cd in $PATH）
# 说明 cd 不是独立程序，而是 Shell 内建命令
\end{lstlisting}

\textbf{同类内建命令}：\code{cd}、\code{export}、\code{alias}、\code{exit}、\code{source}/\code{.} 等——凡是需要修改 Shell 自身状态的命令，都必须是内建命令。
\end{exercise}

\begin{exercise}{写脚本检查文件是否存在}
\textbf{要求}：写一个脚本，接收文件名参数（\code{\$1}），用 \code{test -f} 或 \code{[ -f ... ]} 检查该文件是否存在，并根据结果输出不同提示。

\textbf{脚本 \code{check.sh}}：
\begin{lstlisting}
#!/bin/bash
# 检查文件是否存在的脚本

if [ -f "$1" ]; then
    echo "文件 '$1' 存在。"
else
    echo "文件 '$1' 不存在。"
fi
\end{lstlisting}

\textbf{关键点}：
\begin{itemize}[leftmargin=2em]
    \item \code{\$1} 是第一个命令行参数
    \item \code{[ -f "\$1" ]} 测试 \code{\$1} 是否为存在的\textbf{普通文件}（\code{-f} = file）
    \item \code{"\$1"} 加双引号防止文件名含空格时出错
    \item \code{if ...; then ... else ... fi} 是 bash 条件语句结构
\end{itemize}

\textbf{其他常用测试条件}：
\begin{table}[h]
\centering
\begin{tabular}{c|l}
\toprule
\textbf{条件} & \textbf{含义} \\
\midrule
\code{-f file} & 存在且是普通文件 \\
\code{-d file} & 存在且是目录 \\
\code{-e file} & 存在（任何类型） \\
\code{-r file} & 存在且可读 \\
\code{-w file} & 存在且可写 \\
\code{-x file} & 存在且可执行 \\
\code{-s file} & 存在且大小大于0 \\
\bottomrule
\end{tabular}
\end{table}
\end{exercise}

\begin{exercise}{\code{chmod +x} 的作用}
\textbf{要求}：把上一题的脚本保存为 \code{check.sh}。先运行 \code{./check.sh somefile}，会发生什么？然后执行 \code{chmod +x check.sh} 再试一次。为什么这一步是必须的？

\textbf{chmod 前运行}：
\begin{lstlisting}
./check.sh somefile
# 输出: bash: ./check.sh: Permission denied
\end{lstlisting}

\textbf{chmod 前的权限}：
\begin{lstlisting}
-rw-------  1 liuchenxi  staff  147 Aug 30 16:52 check.sh
\end{lstlisting}
\code{rw-} = 所有者可读、可写，但\textbf{不可执行}

\textbf{执行 \code{chmod +x check.sh} 后}：
\begin{lstlisting}
-rwx--x--x  1 liuchenxi  staff  147 Aug 30 16:52 check.sh
\end{lstlisting}
\code{rwx} = 所有者可读、可写、可执行

\textbf{chmod 后运行}：
\begin{lstlisting}
./check.sh somefile
# 输出: 文件 'somefile' 不存在。

./check.sh check.sh
# 输出: 文件 'check.sh' 存在。
\end{lstlisting}

\textbf{为什么必须 \code{chmod +x}}：
\begin{itemize}[leftmargin=2em]
    \item Unix 系统中，文件能否被直接执行取决于\textbf{可执行权限位}（x 位），而不是文件扩展名
    \item 新建文件默认没有执行权限（安全考虑）
    \item \code{chmod +x} 给所有用户添加执行权限，Shell 才能通过 \code{./check.sh} 直接运行它
    \item 没有执行权限时，仍可通过解释器显式运行：\code{bash check.sh somefile}
\end{itemize}

\noindent\textit{shebang 行（\code{\#!/bin/bash}）只在文件有执行权限且被直接运行时才生效。}
\end{exercise}

\begin{exercise}{\code{set -x} 的作用}
\textbf{要求}：在脚本的 \code{set} 选项里加入 \code{-x} 会发生什么？写个简单脚本试试。

\textbf{\code{set -x} 的作用}：开启\textbf{调试模式}（xtrace）。Shell 在执行每条命令前，会先把该命令及其展开后的参数打印到 stderr（通常以 \code{+} 开头），方便追踪脚本执行流程。

\textbf{演示脚本 \code{debug\_demo.sh}}：
\begin{lstlisting}
#!/bin/bash
set -x

echo "开始执行脚本"
name="Shell"
echo "你好，$name"
for i in 1 2 3; do
    echo "第 $i 次循环"
done
echo "脚本结束"
\end{lstlisting}

\textbf{运行输出}（\code{+} 开头的行是 set -x 打印的追踪信息）：
\begin{lstlisting}
+ echo '开始执行脚本'
开始执行脚本
+ name=Shell
+ echo '你好，Shell'
你好，Shell
+ i=1
+ echo '第 1 次循环'
第 1 次循环
+ i=2
+ echo '第 2 次循环'
第 2 次循环
+ i=3
+ echo '第 3 次循环'
第 3 次循环
+ echo '脚本结束'
脚本结束
\end{lstlisting}

\textbf{常用 set 选项总结}：
\begin{table}[h]
\centering
\begin{tabular}{c|l}
\toprule
\textbf{选项} & \textbf{作用} \\
\midrule
\code{set -e} & 任何命令失败时立即退出脚本 \\
\code{set -u} & 使用未定义变量时报错（而非当作空字符串） \\
\code{set -o pipefail} & 管道中任一命令失败，整个管道视为失败 \\
\code{set -x} & 打印每条执行的命令（调试用） \\
\bottomrule
\end{tabular}
\end{table}

\noindent\textbf{推荐组合}：\code{set -euo pipefail}（脚本开头的「严格模式」），调试时再加 \code{-x}。
\end{exercise}

\begin{exercise}{带日期的备份文件}
\textbf{要求}：写一条命令，把文件复制为带当天日期的备份文件名（例如 \code{notes.txt} $\rightarrow$ \code{notes\_2026-01-12.txt}）。提示：\code{\$(date +\%Y-\%m-\%d)}

\textbf{命令}：
\begin{lstlisting}
cp notes.txt "notes_$(date +%Y-%m-%d).txt"
\end{lstlisting}

\textbf{原理}：
\begin{itemize}[leftmargin=2em]
    \item \code{\$(date +\%Y-\%m-\%d)} 是\textbf{命令替换}：执行 \code{date} 命令，用其输出替换整个 \code{\$()}
    \item \code{+\%Y-\%m-\%d} 是 date 的格式参数：\code{\%Y}=四位年, \code{\%m}=两位月, \code{\%d}=两位日
    \item 双引号确保日期展开后作为一个整体参数
\end{itemize}

\textbf{实际运行}：
\begin{lstlisting}
echo "这是笔记内容" > notes.txt
cp notes.txt "notes_$(date +%Y-%m-%d).txt"
ls -l notes*.txt
\end{lstlisting}

\textbf{输出}：
\begin{lstlisting}
-rw-r--r--  1 staff  19 Aug 30 16:52 notes.txt
-rw-r--r--  1 staff  19 Aug 30 16:52 notes_2026-08-30.txt
\end{lstlisting}

\textbf{通用公式}：
\begin{lstlisting}
cp "文件名" "文件名_$(date +%Y-%m-%d).扩展名"
\end{lstlisting}

\textbf{其他常用日期格式}：
\begin{lstlisting}
date +%Y%m%d                  # 20260830
date +%Y-%m-%d_%H%M%S         # 2026-08-30_165200（带时间）
date -v-1d +%Y-%m-%d          # 昨天的日期（macOS）
\end{lstlisting}
\end{exercise}

\begin{exercise}{修改 flaky test 脚本接收命令行参数}
\textbf{要求}：修改讲义中的「复现偶尔才会失败的测试」脚本，使它能够从命令行参数接收测试命令，而不是写死 \code{cargo test my\_test}。提示：\code{\$1} 或 \code{\$@}

\textbf{修改后的脚本 \code{flaky\_test.sh}}：
\begin{lstlisting}
#!/bin/bash
set -euo pipefail

# 从命令行参数接收测试命令，而不是写死
TEST_CMD="$@"

if [ -z "$TEST_CMD" ]; then
    echo "用法: $0 <测试命令>"
    echo "示例: $0 cargo test my_test"
    exit 1
fi

# Start CPU stress in background
stress --cpu 8 &
STRESS_PID=$!

# Setup log file
LOGFILE="test_runs_$(date +%s).log"
echo "Logging to $LOGFILE"
echo "运行测试命令: $TEST_CMD"

# Run tests until one fails
RUN=1
while $TEST_CMD > "$LOGFILE" 2>&1; do
    echo "Run $RUN passed"
    ((RUN++))
done

# Cleanup and report
kill $STRESS_PID
echo "Test failed on run $RUN"
echo "Last 20 lines of output:"
tail -n 20 "$LOGFILE"
echo "Full log: $LOGFILE"
\end{lstlisting}

\textbf{关键改动}：
\begin{itemize}[leftmargin=2em]
    \item 原脚本写死：\code{while cargo test my\_test > "\$LOGFILE" 2>\&1; do}
    \item 改为：\code{TEST\_CMD="\$@"} 接收所有命令行参数，然后 \code{while \$TEST\_CMD > ...}
\end{itemize}

\textbf{特殊参数说明}：
\begin{table}[h]
\centering
\begin{tabular}{c|l}
\toprule
\textbf{参数} & \textbf{含义} \\
\midrule
\code{\$1} & 第一个参数 \\
\code{\$@} & 所有参数，每个参数独立（加引号后 \code{"\$@"} 保留空格） \\
\code{\$*} & 所有参数，用 IFS 连接成一个字符串 \\
\code{\$#} & 参数个数 \\
\code{\$0} & 脚本自身名称 \\
\code{\$!} & 最近一个后台进程的 PID \\
\bottomrule
\end{tabular}
\end{table}

\textbf{使用方式}：
\begin{lstlisting}
./flaky_test.sh cargo test my_test     # 传 cargo test 命令
./flaky_test.sh pytest tests/           # 传 pytest 命令
./flaky_test.sh false                   # 用 false 测试（第一次就失败）
\end{lstlisting}

\noindent\textit{注意：\code{\$@} 比 \code{\$1} 更适合接收完整命令，因为命令可能包含多个参数（如 \code{cargo test my\_test} 是三个参数）。}
\end{exercise}

\section{管道、工具与进阶}

\begin{exercise}{找出 home 目录最常见的 5 种文件扩展名}
\textbf{要求}：使用管道找出 home 目录中最常见的 5 种文件扩展名。提示：组合 \code{find}、\code{grep}/\code{sed}/\code{awk}、\code{sort}、\code{uniq -c}、\code{head}

\textbf{命令}：
\begin{lstlisting}
find ~ -type f 2>/dev/null \
  | grep -oE '\.[^./]+$' \
  | sort \
  | uniq -c \
  | sort -rn \
  | head -n 5
\end{lstlisting}

\textbf{逐步拆解}：
\begin{enumerate}[leftmargin=2em]
    \item \code{find ~ -type f}：递归列出 home 目录下所有普通文件
    \item \code{2>/dev/null}：把权限错误等 stderr 丢弃
    \item \code{grep -oE '\textbackslash .[\^{}./]+\$'}：提取每行末尾的扩展名
    \begin{itemize}[leftmargin=1.5em]
        \item \code{-o}：只输出匹配部分
        \item \code{-E}：使用扩展正则
        \item \code{\textbackslash .[\^{}./]+\$}：匹配 \code{.} + 一个或多个非\code{.}非\code{/}字符 + 行尾
    \end{itemize}
    \item \code{sort}：排序（相同扩展名排在一起）
    \item \code{uniq -c}：去重并统计每行出现次数
    \item \code{sort -rn}：按数字逆序排序（出现次数多的在前）
    \item \code{head -n 5}：取前 5 个
\end{enumerate}

\textbf{实际输出}：
\begin{lstlisting}
11425 .py
 9438 .js
 6911 .pyc
 5855 .md
 3948 .json
\end{lstlisting}

\textbf{等价的 awk 版本}：
\begin{lstlisting}
find ~ -type f 2>/dev/null | awk -F. '{print "."$NF}' | sort | uniq -c | sort -rn | head -5
\end{lstlisting}
\end{exercise}

\begin{exercise}{\code{xargs} 结合 \code{find} 统计 \code{.sh} 文件行数}
\textbf{要求}：\code{xargs} 把 stdin 的每一行转换为命令参数。结合 \code{find} 和 \code{xargs}（不要用 \code{find -exec}），找出目录中所有 \code{.sh} 文件，并用 \code{wc -l} 统计每个文件行数。加分项：正确处理文件名中的空格。提示：\code{-print0} 和 \code{-0}

\textbf{基础版本}：
\begin{lstlisting}
find . -name "*.sh" | xargs wc -l
\end{lstlisting}

\textbf{问题}：如果文件名包含空格（如 \code{my script.sh}），xargs 会把它拆成两个参数 \code{my} 和 \code{script.sh}，导致错误。

\textbf{正确处理空格的版本（加分项）}：
\begin{lstlisting}
find . -name "*.sh" -print0 | xargs -0 wc -l
\end{lstlisting}

\textbf{原理}：
\begin{itemize}[leftmargin=2em]
    \item \code{find -print0}：用\textbf{null 字符（\textbackslash0）}分隔文件名，而不是换行符
    \item \code{xargs -0}：以 null 字符作为分隔符读取输入
    \item 因为文件名中不可能包含 null 字符，所以这种方式\textbf{100\% 安全}处理任何文件名（含空格、换行、特殊字符）
\end{itemize}

\textbf{实际输出}：
\begin{lstlisting}
      1 ./dir with spaces/my script.sh
     34 ./flaky_test.sh
      8 ./check.sh
     10 ./debug_demo.sh
     53 total
\end{lstlisting}

\textbf{xargs 常用参数}：
\begin{table}[h]
\centering
\begin{tabular}{c|l}
\toprule
\textbf{参数} & \textbf{含义} \\
\midrule
\code{-0} & 以 null 字符分隔输入（配合 find -print0） \\
\code{-n N} & 每次传递 N 个参数给命令 \\
\code{-I \{\}} & 用 \code{\{\}} 作为参数占位符 \\
\code{-P N} & 并行运行 N 个进程 \\
\bottomrule
\end{tabular}
\end{table}
\end{exercise}

\begin{exercise}{\code{curl} 获取课程网站 HTML，统计列出了多少讲}
\textbf{要求}：使用 \code{curl} 获取课程网站的 HTML，通过 \code{grep} 统计列出了多少讲。提示：找出每讲课程名称在 HTML 中的共性；用 \code{curl -s} 关闭进度输出。

\textbf{命令}：
\begin{lstlisting}
curl -s https://missing-semester-cn.github.io/2026/ \
  | grep -oE 'href="/2026/[a-z-]+/"' \
  | sort -u \
  | wc -l
\end{lstlisting}

\textbf{逐步拆解}：
\begin{enumerate}[leftmargin=2em]
    \item \code{curl -s URL}：静默模式获取网页 HTML（不显示进度条）
    \item \code{grep -oE 'href="/2026/[a-z-]+/"'}：提取所有课程链接
    \begin{itemize}[leftmargin=1.5em]
        \item 课程链接的共性：\code{href="/2026/课程名/"}
        \item 课程名由小写字母和连字符组成
    \end{itemize}
    \item \code{sort -u}：排序并去重（\code{-u} = unique）
    \item \code{wc -l}：统计行数，即课程讲数
\end{enumerate}

\textbf{实际输出}：
\begin{lstlisting}
9
\end{lstlisting}

\textbf{列出各讲名称}：
\begin{lstlisting}
curl -s https://missing-semester-cn.github.io/2026/ \
  | grep -oE 'href="/2026/[a-z-]+/"' \
  | sort -u \
  | sed 's|href="/2026/||;s|/"||'
\end{lstlisting}

\textbf{输出}：
\begin{lstlisting}
agentic-coding
beyond-code
code-quality
command-line-environment
course-shell
debugging-profiling
development-environment
shipping-code
version-control
\end{lstlisting}

共 \textbf{9 讲}，与课程介绍中「九场 1 小时的讲座」一致。
\end{exercise}

\begin{exercise}{\code{jq} 处理 JSON}
\textbf{要求}：用 curl 获取示例数据，再用 jq 提取 version 大于 6 的人员姓名。提示：先 \code{jq .} 看结构；再试 \code{jq '.[] | select(...) | .name'}

\textbf{获取数据并查看结构}：
\begin{lstlisting}
curl -s https://microsoftedge.github.io/Demos/json-dummy-data/64KB.json > data.json
jq '.[0]' data.json
\end{lstlisting}

\textbf{JSON 结构}：
\begin{lstlisting}
{
  "name": "Adeel Solangi",
  "language": "Sindhi",
  "id": "V59OF92YF627HFY0",
  "bio": "Donec lobortis eleifend...",
  "version": 6.1
}
\end{lstlisting}
顶层是一个数组，每个元素包含 name、language、id、bio、version 字段。

\textbf{提取 version > 6 的人员姓名}：
\begin{lstlisting}
jq '.[] | select(.version > 6) | .name' data.json
\end{lstlisting}

\textbf{jq 语法拆解}：
\begin{itemize}[leftmargin=2em]
    \item \code{.[]}：遍历数组中的每个元素
    \item \code{select(.version > 6)}：筛选 version 字段大于 6 的元素
    \item \code{.name}：提取 name 字段
\end{itemize}

\textbf{统计数量}：
\begin{lstlisting}
jq '[.[] | select(.version > 6)] | length' data.json
# 输出: 89
\end{lstlisting}

\textbf{部分输出}：
\begin{lstlisting}
"Adeel Solangi"
"Aamir Solangi"
"Adil Eli"
"George Mifsud"
"John Falzon"
...
\end{lstlisting}

\textbf{jq 常用操作速查}：
\begin{table}[h]
\centering
\begin{tabular}{c|l}
\toprule
\textbf{表达式} & \textbf{含义} \\
\midrule
\code{.} & 整个输入 \\
\code{.field} & 取字段 \\
\code{.[]} & 遍历数组 \\
\code{select(cond)} & 筛选满足条件的元素 \\
\code{length} & 数组长度或字符串长度 \\
\code{keys} & 对象的所有键 \\
\code{map(expr)} & 对数组每个元素应用 expr \\
\bottomrule
\end{tabular}
\end{table}
\end{exercise}

\begin{exercise}{\code{awk} 按列过滤并交换列}
\textbf{要求}：写一个 awk 命令：只输出第二列大于 100 的行，并交换第一列和第三列。测试数据：\code{printf 'a 50 x\textbackslash nb 150 y\textbackslash nc 200 z\textbackslash n'}

\textbf{命令}：
\begin{lstlisting}
printf 'a 50 x\nb 150 y\nc 200 z\n' \
  | awk '$2 > 100 { temp=$1; $1=$3; $3=temp; print }'
\end{lstlisting}

\textbf{awk 语法拆解}：
\begin{itemize}[leftmargin=2em]
    \item \code{\$2 > 100}：模式（条件），只处理第二列大于 100 的行
    \item \code{\{ temp=\$1; \$1=\$3; \$3=temp; print \}}：动作
    \begin{itemize}[leftmargin=1.5em]
        \item \code{temp=\$1}：把第一列存到临时变量
        \item \code{\$1=\$3}：第一列替换为第三列的值
        \item \code{\$3=temp}：第三列替换为原来的第一列
        \item \code{print}：打印修改后的整行
    \end{itemize}
\end{itemize}

\textbf{实际输出}：
\begin{lstlisting}
y 150 b
z 200 c
\end{lstlisting}

\textbf{验证}：
\begin{itemize}[leftmargin=2em]
    \item 原始行 \code{a 50 x}：第二列 50，不大于 100，被过滤
    \item 原始行 \code{b 150 y}：第二列 150 > 100，交换第一列(b)和第三列(y) $\rightarrow$ \code{y 150 b}
    \item 原始行 \code{c 200 z}：第二列 200 > 100，交换第一列(c)和第三列(z) $\rightarrow$ \code{z 200 c}
\end{itemize}

\textbf{awk 核心概念}：
\begin{itemize}[leftmargin=2em]
    \item awk 按行处理，每行自动按空白字符分成列
    \item \code{\$1, \$2, \$3...} 引用各列，\code{\$0} 表示整行
    \item \code{NR} = 当前行号，\code{NF} = 当前行列数
    \item \code{FS} = 输入分隔符（默认空白），\code{-F,} 可设为逗号
    \item \code{BEGIN \{\}} 在处理前执行，\code{END \{\}} 在处理后执行
\end{itemize}
\end{exercise}

\begin{exercise}{拆解 SSH 日志管道 + 从历史记录找出最常用命令}
\textbf{要求}：拆解讲义中的 SSH 日志处理管道：每一步分别做了什么？然后仿照它构建一个管道，从 \code{\textasciitilde/.bash\_history}（或 \code{\textasciitilde/.zsh\_history}）中找出最常使用的 Shell 命令。

\subsection*{第一部分：拆解 SSH 日志管道}

\textbf{原始管道}：
\begin{lstlisting}
ssh myserver 'journalctl -u sshd -b-1 | grep "Disconnected from"' \
  | sed -E 's/.*Disconnected from .* user (.*) [^ ]+ port.*/\1/' \
  | sort | uniq -c \
  | sort -nk1,1 | tail -n10 \
  | awk '{print $2}' | paste -sd,
\end{lstlisting}

\textbf{逐步拆解}：
\begin{table}[h]
\centering
\begin{tabular}{c|l|l}
\toprule
\textbf{步骤} & \textbf{命令} & \textbf{作用} \\
\midrule
1 & \code{ssh myserver '...'} & 登录远程服务器，在远程执行引号内的命令 \\
2 & \code{journalctl -u sshd -b-1} & 查看 sshd 服务在上一次启动时的系统日志 \\
3 & \code{grep "Disconnected from"} & 筛选包含「断开连接」的日志行 \\
4 & \code{sed -E 's/...user (.*) .../\1/'} & 用正则提取用户名，\code{(.*)} 捕获，\code{\textbackslash1} 输出 \\
5 & \code{sort} & 对用户名排序（相同名字排在一起） \\
6 & \code{uniq -c} & 去重并统计每个用户名出现次数 \\
7 & \code{sort -nk1,1} & 按第一列（次数）数值排序（升序） \\
8 & \code{tail -n10} & 取最后 10 个，即出现次数最多的 10 个 \\
9 & \code{awk '\{print \$2\}'} & 只输出第二列（用户名），去掉次数 \\
10 & \code{paste -sd,} & 把多行合并为一行，用逗号分隔 \\
\bottomrule
\end{tabular}
\end{table}

\textbf{最终输出示例}：
\begin{lstlisting}
postgres,mysql,oracle,dell,ubuntu,inspur,test,admin,user,root
\end{lstlisting}

\subsection*{第二部分：从历史记录找出最常用命令}

\textbf{命令（zsh 用户）}：
\begin{lstlisting}
LC_ALL=C awk '{print $1}' ~/.zsh_history | sort | uniq -c | sort -rn | head -n 10
\end{lstlisting}

\textbf{命令（bash 用户）}：
\begin{lstlisting}
awk '{print $1}' ~/.bash_history | sort | uniq -c | sort -rn | head -n 10
\end{lstlisting}

\textbf{逐步拆解}：
\begin{enumerate}[leftmargin=2em]
    \item \code{awk '\{print \$1\}'}：提取每行的第一个词（即命令名，去掉参数）
    \item \code{sort}：排序
    \item \code{uniq -c}：统计每个命令出现次数
    \item \code{sort -rn}：按次数逆序排序
    \item \code{head -n 10}：取前 10 个
\end{enumerate}

\textbf{实际输出（zsh\_history）}：
\begin{lstlisting}
123 brew
 25 cd
 22 git
 20 sudo
 14 ls
  9 echo
  8 fc-match
  7 which
  7 find
  6 mkdir
\end{lstlisting}

\textbf{说明}：
\begin{itemize}[leftmargin=2em]
    \item zsh 的历史文件默认格式可能包含时间戳前缀（\code{: 1234567890:0;command}），如果有，需要先用 \code{sed 's/\^{}: [0-9]*:[0-9]*;//'} 去掉前缀
    \item \code{LC\_ALL=C} 防止遇到非 UTF-8 字节时 awk/sed 报错
    \item 最常用的命令是 \code{brew}（macOS 包管理器），其次是 \code{cd}、\code{git}、\code{sudo}、\code{ls}
\end{itemize}
\end{exercise}

\section*{附录：本练习涉及的脚本文件}
\addcontentsline{toc}{section}{附录：脚本文件}

\begin{table}[h]
\centering
\begin{tabular}{c|l}
\toprule
\textbf{文件} & \textbf{说明} \\
\midrule
\code{check.sh} & 题8/9：检查文件是否存在的脚本 \\
\code{debug\_demo.sh} & 题10：演示 set -x 调试模式 \\
\code{flaky\_test.sh} & 题12：接收命令行参数的 flaky test 脚本 \\
\bottomrule
\end{tabular}
\end{table}

所有脚本均已添加可执行权限，可直接运行。

\vspace{2em}
\begin{center}
\textit{--- 全文完 ---}
\end{center}

\end{document}
